\section{多层 RNN}

\begin{figure}[H]
\centering
\includegraphics[width=0.95\textwidth]{slides-images/slide-046.jpg}
\caption{多层 RNN 结构：每一层有自己的权重矩阵，层内时间步共享权重，层间通过输出传递信息\protect\footnotemark}
\end{figure}
\footnotetext{Slides 第47页。}

多层 RNN 的组织方式如下：
\begin{itemize}
    \item \textbf{层内}：同一层的所有时间步共享权重，隐藏状态沿时间维度传递
    \item \textbf{层间}：第 $l$ 层的输出作为第 $l+1$ 层的输入
    \item 每一层有\textbf{独立}的权重矩阵
\end{itemize}

这形成了一个二维的计算网格（时间 $\times$ 深度），进一步加剧了训练时的计算瓶颈。

\subsection{本章小结}

多层 RNN 通过堆叠 RNN 层来增加模型的表达能力，但也使得计算图更加复杂，训练更加困难。

%% ========== 梯度问题与 LSTM ==========

